\appendix
\section{Proofs and Numerical Validation}
\label{app:proofs}

\subsection{Proof of \cref{prop:expansion}}
Apply Taylor's theorem with Lagrange remainder to
$t\mapsto\loss(f_A(x)+t\,\dAj(x),y)$ on $[0,1]$. There exists
$\xi=\xi(x,y)$ on the segment with
\begin{equation}
 \loss(f_{A\cup\{j\}},y)
 =
 \loss(f_A,y)+\grad\loss(f_A,y)^\top \dAj
 +\tfrac12 \dAj^\top\grad^2\loss(\xi,y)\dAj .
\end{equation}
Rearranging, subtracting $\lambda c_j$ and taking expectations gives
\eqref{eq:prop-expansion}. Under the $L$-smoothness assumption
$|\dAj^\top\grad^2\loss(\xi,y)\dAj|\leq L\|\dAj\|^2$ for every realisation,
which yields \eqref{eq:prop-bound}.

\subsection{Proof of \cref{prop:bregman}}
The Bregman three-point identity states that for any $u,v,w$,
$D_\phi(u,v)-D_\phi(u,w)=(\grad\phi(w)-\grad\phi(v))^\top(u-v)-D_\phi(v,w)$;
both sides expand to the difference of
$\phi(u)-\phi(v)-\langle\grad\phi(v),u-v\rangle$ and
$\phi(u)-\phi(w)-\langle\grad\phi(w),u-w\rangle$. Set $u=y$, $v=f_A$ and
$w=f_{A\cup\{j\}}$, take expectations, and subtract $\lambda c_j$.

\subsection{Proof of \cref{prop:ce}}
For $\loss(f,y)=-f_y+\log\sum_k e^{f_k}$ we have $\grad\loss=p-e_y$ and
$\grad^2\loss=\diag(p)-pp^\top\succeq0$. Since
$(\diag(p)-pp^\top)\mathbf 1=0$, the quadratic form is invariant to adding a
constant to $v$, so $v^\top(\diag(p)-pp^\top)v=\Var_p(v)$ for the random
variable taking value $v_k$ with probability $p_k$. Let $M=\max_k v_k$,
$m=\min_k v_k$ and $\mu_p=\sum_k p_k v_k$. Bhatia--Davis gives
$\Var_p(v)\leq(M-\mu_p)(\mu_p-m)\leq(M-m)^2/4$. Replacing $v$ by
$v-\bar v\mathbf 1$ leaves the quadratic form and the range bound valid and
makes the vector zero-sum, so $M\geq0\geq m$, $\|v\|^2\geq M^2+m^2$, and
$(M-m)^2=(M+|m|)^2\leq2(M^2+m^2)\leq2\|v\|^2$. Hence
$v^\top(\diag(p)-pp^\top)v\leq\tfrac12\|v\|^2$; the bound is attained at
$p=(\tfrac12,\tfrac12,0,\dots)$, $v=(1,-1,0,\dots)/\sqrt2$, so the constant is
tight. Applying \cref{prop:expansion} with $L=\tfrac12$ and
$-\grad\loss(f_A,y)=e_y-p_A$ gives \eqref{eq:prop-ce}.

\subsection{Proof of \cref{prop:suff} and \cref{cor:identifiability}}
For squared loss \eqref{eq:prop-bregman} reads
$\widetilde\marginal=\loss(f_A,y)-\loss(f_{A\cup\{j\}},y)-\lambda c_j
=2\dAj^\top(y-f_A)-\|\dAj\|^2-\lambda c_j$. The response $\dAj=\dAj(x)$ is
measurable with respect to $(x,A)$, so it can be taken outside the conditional
expectation; the tower property gives \eqref{eq:prop-suff} with
$\mu(x)=\E[y\mid x,A]$. Linearity in $\dAj$ and the response-only quadratic
penalty follow immediately. If the predictor is conditionally calibrated,
$\mu(x)=f_A(x)$ for every observable state, so $g_A\equiv0$ and the expected
utility reduces to $-\|\dAj\|^2-\lambda c_j$.

\subsection{Numerical validation}
\label{app:validation}

\paragraph{Squared loss (exactness).} Using the subset-capable nonlinear
forecasting expert of \cref{sec:experiments} on held-out episodes, we compared
the exact marginal $\loss(f_A,y)-\loss(f_{A\cup\{j\}},y)$ with the right-hand
side of the squared-loss identity of \cref{prop:bregman}. Over 384 held-out
(state, candidate) pairs the maximum absolute deviation is $3.3\times10^{-7}$ in
normalised units. Unit tests enforce the same identity, the agreement between
batched pair evaluation and explicit augmented-set evaluation, and the
inertness of unselected candidates.

\paragraph{Cross-entropy (two-sided bound).} For the frozen in-context
classifier used in the demonstration-selection experiments we collected 92{,}995 (state, candidate) samples together
with 2{,}975{,}840 per-query rows and evaluated the sandwich of
\cref{prop:ce}. The bound held in every sample at both aggregation levels
(smallest upper slack $3.0\times10^{-8}$, smallest lower slack
$1.0\times10^{-3}$), and the first-order term tracks the true marginal closely:
Spearman $.971$ and $R^2=.853$ pooled, with per-benchmark Spearman between
$.970$ and $.993$.

\paragraph{Hessian constant.} The spectral bound of \cref{prop:ce} was checked
numerically by maximising the Rayleigh quotient over the simplex for
$C\in\{2,3,5,10,50\}$; the maximum is $0.5$ in every case, confirming that the
constant is tight.

\paragraph{Response sufficiency.} For the squared-loss expert we fitted the same
ridge head twice on identical held-out state samples, once with state features
only and once with state plus response features. The response features raise the
out-of-sample $R^2$ from $-.89$ to $-.22$ on a held-out episode split, a gain of
$.67$; the cross-entropy family shows the same direction with a Spearman
correlation of $.97$ between the first-order term and the realised marginal.

\begin{table}[t]
 \centering
 \small
 \caption{Numerical validation of the response--utility relations.}
 \label{tab:theory-validation-app}
 \begin{tabular}{lll}
  \toprule
  Check & Object & Result \\
  \midrule
  Squared-loss identity & forecast expert & max.\ deviation $3.3\times10^{-7}$\\
  Cross-entropy sandwich & in-context classifier & 0 violations / 92{,}995\\
  First-order vs.\ marginal & in-context classifier & Spearman .971, $R^2$ .853\\
  Response vs.\ state only & forecast expert & $R^2$ gain $+0.67$\\
  Hessian spectral bound & softmax, all $C\le 50$ & maximum $=0.5$\\
  \bottomrule
 \end{tabular}
\end{table}


\section{Additional Stress Tests}
\label{app:stress-tests}

The main text keeps the predictor-strength story compact. This section records
the additional tests used to eliminate alternative explanations.

\paragraph{Search direction.}
Backward elimination from the full pool and forward greedy selection differ by
only $+.0002$ on the matched cells, with the interval spanning zero. The search
direction has little effect once the utility signal itself is fixed.

\paragraph{Harmful candidate pools.}
Rebuilding the candidate pools from weakly or negatively correlated sensors does
not restore a reliable routing advantage under the strong predictor. Pooling
remains competitive, and the response-aware selector stays far below the oracle.
The result shows that robustness to obviously poor candidates does not explain
the missing realizable value.

\paragraph{Response-geometry limit case.}
When the predictable residual shrinks, \cref{prop:suff} leaves the quadratic
response term as the observable part of the score. Selecting the least
perturbing candidates follows that limit directly. Its gain is $+.0026$ on
METR-LA and $+.0037$ on PEMS-BAY, both below pooling. Response magnitude alone
preserves the sign of the curvature effect but does not identify which candidate
adds useful information.

\paragraph{Full-pool reranking.}
Setting $q=K$ gives the response head access to every candidate. The gain changes
by $+.0012$ with a confidence interval that includes zero relative to the
shortlist version. The shortlist is therefore an efficiency device, while the
ranking ceiling is set by the observable signal.

\paragraph{Episode-level confidence.}
We record the cached top-two margin, response-aware top-two margin, best utility
score and prediction-change norm for every episode. Across datasets and signals,
AUCs for predicting whether routing beats pooling lie in $[.481,.521]$. A
held-out confidence rule recovers only a small fraction of the oracle
episode-level switching gain. The selector's own confidence does not reveal its
mistakes once the predictor has absorbed the systematic residual.

\paragraph{Policy-gradient alternative.}
A REINFORCE selector trained directly on subset reward performs below random
selection in the same strong-predictor setting. This rules out the possibility
that the observed ceiling comes only from marginal-utility regression.
